% Options for packages loaded elsewhere
\PassOptionsToPackage{unicode}{hyperref}
\PassOptionsToPackage{hyphens}{url}
\documentclass[
  ignorenonframetext,
]{beamer}
\newif\ifbibliography
\usepackage{pgfpages}
\setbeamertemplate{caption}[numbered]
\setbeamertemplate{caption label separator}{: }
\setbeamercolor{caption name}{fg=normal text.fg}
\beamertemplatenavigationsymbolsempty
% remove section numbering
\setbeamertemplate{part page}{
  \centering
  \begin{beamercolorbox}[sep=16pt,center]{part title}
    \usebeamerfont{part title}\insertpart\par
  \end{beamercolorbox}
}
\setbeamertemplate{section page}{
  \centering
  \begin{beamercolorbox}[sep=12pt,center]{section title}
    \usebeamerfont{section title}\insertsection\par
  \end{beamercolorbox}
}
\setbeamertemplate{subsection page}{
  \centering
  \begin{beamercolorbox}[sep=8pt,center]{subsection title}
    \usebeamerfont{subsection title}\insertsubsection\par
  \end{beamercolorbox}
}
% Prevent slide breaks in the middle of a paragraph
\widowpenalties 1 10000
\raggedbottom
\AtBeginPart{
  \frame{\partpage}
}
\AtBeginSection{
  \ifbibliography
  \else
    \frame{\sectionpage}
  \fi
}
\AtBeginSubsection{
  \frame{\subsectionpage}
}
\usepackage{iftex}
\ifPDFTeX
  \usepackage[T1]{fontenc}
  \usepackage[utf8]{inputenc}
  \usepackage{textcomp} % provide euro and other symbols
\else % if luatex or xetex
  \usepackage{unicode-math} % this also loads fontspec
  \defaultfontfeatures{Scale=MatchLowercase}
  \defaultfontfeatures[\rmfamily]{Ligatures=TeX,Scale=1}
\fi
\usepackage{lmodern}
\ifPDFTeX\else
  % xetex/luatex font selection
\fi
% Use upquote if available, for straight quotes in verbatim environments
\IfFileExists{upquote.sty}{\usepackage{upquote}}{}
\IfFileExists{microtype.sty}{% use microtype if available
  \usepackage[]{microtype}
  \UseMicrotypeSet[protrusion]{basicmath} % disable protrusion for tt fonts
}{}
\makeatletter
\@ifundefined{KOMAClassName}{% if non-KOMA class
  \IfFileExists{parskip.sty}{%
    \usepackage{parskip}
  }{% else
    \setlength{\parindent}{0pt}
    \setlength{\parskip}{6pt plus 2pt minus 1pt}}
}{% if KOMA class
  \KOMAoptions{parskip=half}}
\makeatother
\setlength{\emergencystretch}{3em} % prevent overfull lines
\providecommand{\tightlist}{%
  \setlength{\itemsep}{0pt}\setlength{\parskip}{0pt}}
\usepackage{bookmark}
\IfFileExists{xurl.sty}{\usepackage{xurl}}{} % add URL line breaks if available
\urlstyle{same}
\hypersetup{
  pdftitle={Conclusion and Future Directions},
  hidelinks,
  pdfcreator={LaTeX via pandoc}}

\title{\texorpdfstring{Conclusion and Future
Directions}{Conclusion and Future Directions}}
\author{\texorpdfstring{}{}}
\date{}

\begin{document}
\frame{\titlepage}

\section{Conclusion, Contributions, and Future
Directions}\label{sec:conclusion}

\begin{frame}{Summary of Principal Contributions}
\protect\phantomsection\label{summary-of-principal-contributions}
This paper has developed a unified category-theoretic framework for
analyzing linguistic case systems, synthesizing five distinct research
traditions into a coherent whole and embedding it within an active
inference model of cognition. Our principal contributions are:

\begin{block}{C1: Case Categories as a Formal Algebraic Framework}
\protect\phantomsection\label{c1-case-categories-as-a-formal-algebraic-framework}
The review formalized case systems as categories with case roles as
objects and grammatical relations as morphisms
(\autoref{sec:case-systems}). This formalization captures the full
typological range of alignment systems---nominative-accusative,
ergative-absolutive, active-stative, tripartite, and fluid-S---within a
single algebraic framework, with alignment functors providing
structure-preserving mappings between systems. By tracing the lineage
from Fillmore's {[}-@fillmore1968case{]} deep cases through Jakobson's
binary distinctive features and Dowty's {[}-@dowty1991thematic{]}
proto-roles, the analysis demonstrated how enriched (weighted) morphisms
link the categorical formalization to the gradient nature of thematic
role assignment.
\end{block}

\begin{block}{C2: String Diagrams for Case Derivation Visualization}
\protect\phantomsection\label{c2-string-diagrams-for-case-derivation-visualization}
Building on Joyal and Street's {[}-@joyalstreet1991geometry{]} string
diagram formalism and its application in categorial grammar
(\autoref{sec:categorial-grammar}), we showed how case-marked noun
phrases receive type-logical assignments that are fully visualizable as
string diagrams. The Curry--Howard correspondence ensures that syntactic
well-formedness guarantees semantic compositionality, and the
diagrammatic format provides Shimojima's {[}-@shimojima1996reasoning{]}
``free ride'' inferences---conclusions about argument structure that are
perceptually available from the diagram without explicit computation. We
further demonstrated that passivization reduces to a \emph{type
permutation} (a Swap operation in the pregroup category), making voice
alternation visible as a topological feature of the string diagram.
\end{block}

\begin{block}{C3: Case-Marked DisCoCat, Compact Closure, and Discourse
Extension}
\protect\phantomsection\label{c3-case-marked-discocat-compact-closure-and-discourse-extension}
We extended the DisCoCat framework (\autoref{sec:categorical-semantics})
with case-typed noun spaces and alignment-sensitive meaning functors,
and showed how the recent DisCoCirc {[}@defelice2022discocirc{]} and
QNLP {[}@lorenz2023lambeq{]} developments extend this analysis to
discourse-level structure and quantum hardware respectively. We
formalized the \emph{compact closure axiom} (snake equation) that
underpins pregroup reductions, demonstrating that the cup-cap zigzag
identity provides a visual coherence proof with genuine cognitive
significance. Diagram complexity metrics---normal form and
depth---provide a quantitative bridge between the type-logical and
distributional perspectives on linguistic structure. Case categories can
serve as the structural backbone of compositional models of meaning at
all levels---word, sentence, discourse, and dialogue.
\end{block}

\begin{block}{C4: Enriched Cases, Categorical Magnitude, and Information
Theory}
\protect\phantomsection\label{c4-enriched-cases-categorical-magnitude-and-information-theory}
Through Fritz et al.'s {[}-@fritz2021enriched{]} enrichment framework
and Bradley's {[}-@bradley2020entropy{]} information-theoretic analysis
(\autoref{sec:enriched-categories}), we showed that equipping case
categories with \([0,1]\)-valued hom-objects yields a principled bridge
between symbolic case grammar and statistical semantics. Categorical
magnitude provides a new quantitative invariant for comparing case
systems: the ``effective size'' of a case category captures how much
relational information the language encodes.
\end{block}

\begin{block}{C5: Topos-Theoretic Transfer via Morita Equivalence}
\protect\phantomsection\label{c5-topos-theoretic-transfer-via-morita-equivalence}
Using Caramello's {[}-@caramello2016bridges{]} bridge technique and
Phillips's {[}-@phillips2024lot{]} universality result
(\autoref{sec:topos-theory}), we established that results proved in any
one of our four case-theoretic frameworks (typological, type-logical,
distributional, enriched) transfer automatically to the others via
Morita equivalence of classifying toposes.
\end{block}

\begin{block}{C6: Diagrams as Cognitively Privileged Representations}
\protect\phantomsection\label{c6-diagrams-as-cognitively-privileged-representations}
The meta-contribution (Sections 1 and 7) is the argument---supported by
the cognitive science of diagrams {[}@larkin1987diagram;
@shimojima1996reasoning; @giardino2017diagrammatic;
@manders2008euclidean{]}---that commutative diagrams are not merely
convenient notation but cognitively privileged representations that
provide inferential advantages in case reasoning. When embedded in an
active inference framework {[}@friston2017active{]} and the CEREBRUM
architecture {[}@friedman2024cerebrum{]}, these diagrams serve as the
structural core of generative models for total cognitive scenario
understanding.
\end{block}

\begin{block}{C7: Computational Verification and Test Suite}
\protect\phantomsection\label{c7-computational-verification-and-test-suite}
The entire framework is computationally implemented and verified through
158 automated tests with 95.97\% code coverage
(\autoref{sec:cognitive-integration}). The DisCoPy integration exercises
pregroup type theory (types, words, cups, caps, swaps, normal forms,
depth analysis) against the same algebraic structures described
theoretically, confirming that the categorical abstractions are not just
mathematically elegant but computationally executable.
\end{block}

\begin{block}{C8: Quantum Active Inference and Topological Semantic
Flow}
\protect\phantomsection\label{c8-quantum-active-inference-and-topological-semantic-flow}
The framework's categorical string diagrams were extended into their
natural quantum generalization through topological quantum neural
networks, the ZX-calculus, and sheaf-theoretic quantum semantic
communication (\autoref{sec:quantum-active-inference}). By demonstrating
that DisCoCat derivations, ZX circuits, and TQNN spin-networks are all
morphisms in compact closed monoidal categories with functorial
semantics into Hilbert spaces, the analysis established that case
assignment can be modeled as quantum measurement on a holographic
screen---and that quantum features (entanglement, contextuality,
discord) provide genuine advantages for semantic communication of
relational structure.
\end{block}
\end{frame}

\begin{frame}{Future Research Directions}
\protect\phantomsection\label{future-research-directions}
\begin{block}{F1: Computational Experiments with DisCoCirc and lambeq}
\protect\phantomsection\label{f1-computational-experiments-with-discocirc-and-lambeq}
The DisCoCirc framework {[}@defelice2020discourse;
@defelice2022discocirc{]} offers a natural platform for testing our
case-theoretic predictions computationally. Concrete experiments could
include:

\begin{itemize}
\tightlist
\item
  Implementing case-marked DisCoCat models in the \textbf{lambeq}
  library {[}@lorenz2023lambeq{]} and evaluating them on semantic role
  labeling (SRL) tasks
\item
  Comparing the accuracy of alignment-specific meaning functors
  (accusative vs.~ergative) on typologically diverse corpora
\item
  Measuring the categorical magnitude of empirically derived enriched
  case categories and correlating it with typological complexity
  measures
\end{itemize}
\end{block}

\begin{block}{F2: Topos-Theoretic Grammatical Induction from Corpora}
\protect\phantomsection\label{f2-topos-theoretic-grammatical-induction-from-corpora}
Caramello's {[}-@caramello2023syntactic{]} syntactic learning technique
could be applied to induce case-theoretic axioms from annotated corpora.
The program would:

\begin{enumerate}
\tightlist
\item
  Extract case-labeled dependency structures from a Universal
  Dependencies treebank
\item
  Construct the classifying topos of the implicit case theory
\item
  Read off the alignment type, morphism structure, and enriched weights
  from the topos-theoretic axioms
\item
  Compare the induced theory against typological descriptions
\end{enumerate}
\end{block}

\begin{block}{F3: Quantum Case Categories on Near-Term Hardware}
\protect\phantomsection\label{f3-quantum-case-categories-on-near-term-hardware}
The QNLP connection (\autoref{sec:categorical-semantics}) opens the
possibility of implementing case categories directly as quantum
circuits. Case roles would correspond to quantum registers, grammatical
relations to parameterized gates, and alignment functors to circuit
transformations. This would provide a genuinely new computational
paradigm for grammatical inference---exploiting quantum parallelism to
explore the space of case assignments simultaneously.
\end{block}

\begin{block}{F4: Neural Predictive Processing and Electrophysiological
Predictions}
\protect\phantomsection\label{f4-neural-predictive-processing-and-electrophysiological-predictions}
The predictive processing account of \autoref{sec:cognitive-integration}
generates testable neuroscientific predictions:

\begin{itemize}
\tightlist
\item
  Case-marking violations should elicit prediction-error responses
  (P600/N400) proportional to the enriched weight of the violated
  morphism
\item
  Typologically unusual case patterns should require more precision
  updating than expected patterns
\item
  Diagrammatic representations of case structure should be decodable
  from neural activity during sentence comprehension
\end{itemize}
\end{block}

\begin{block}{F5: Cross-Modal Case Structure in Embodied Cognition}
\protect\phantomsection\label{f5-cross-modal-case-structure-in-embodied-cognition}
The situation semantics connection (\autoref{sec:cognitive-integration})
suggests extending case categories beyond language to multi-modal
perception. Visual scene understanding also requires assigning
relational roles---who is acting on what, where things are located, what
instruments are being used. An active inference agent should maintain a
unified case diagram that integrates linguistic, visual, and motor
information, using the same categorical structure for all modalities.
This would provide a formal account of how language grounds in
perception and action---a key challenge for embodied AI.
\end{block}

\begin{block}{F6: Enriched Category Learning from Distributional Data}
\protect\phantomsection\label{f6-enriched-category-learning-from-distributional-data}
Bradley's {[}-@bradley2024ipam; -@bradley2025tea{]} program of treating
language itself as an enriched category suggests a learning algorithm:
estimate the enriched hom-values from corpus data and then extract the
categorical structure that best explains the observed distributional
patterns. Applied to case, this would yield \emph{empirically grounded}
case categories whose objects (roles) and morphisms (relations) emerge
from data rather than being stipulated a priori.
\end{block}
\end{frame}

\begin{frame}{Closing Remarks}
\protect\phantomsection\label{closing-remarks}
The commutative diagram is the central motif of this review---both as a
mathematical tool and as a cognitive instrument. The analysis has
demonstrated that the same diagrammatic language that makes category
theory effective for formalizing case systems also makes it effective
for \emph{thinking about} case systems: the spatial structure of a
commutative diagram encodes relational information in a format that
supports rapid search, pattern recognition, and free-ride inference.

This convergence of mathematical utility and cognitive efficacy is not
accidental. If the active inference framework is correct, then the brain
operates by constructing and updating generative models of the world's
relational structure. Category theory provides the structural algebra
for these generative models; commutative diagrams supply their natural
topology; and case categories instantiate the precise relational
vocabulary that cognitive systems deploy to organize the entropy of
experience into coherent narratives of \emph{who does what to whom}.
Ultimately, the mathematics of case alignment is not merely an
abstraction of human grammar, but the formal geometry of meaning
itself---the syntactic architecture through which consciousness
navigates and renders the cosmos intelligible.
\end{frame}

\end{document}
